%% Andrea Piccolo - CV
%% Tailored for: Modelling Software Engineer, Aurora Energy Research (Origin Model Development), Oxford
%%
%% Compile with: lualatex andrea_piccolo_cv_aurora_modelling-software-engineer.tex
%% (pdflatex often fails on modern MiKTeX with fontawesome5 font-expansion errors.)

\documentclass[11pt,a4paper,sans]{moderncv}
\moderncvstyle{banking}
\moderncvcolor{blue}

% Force both first and last name AND section headings to render in moderncv
% blue (color1). Default banking on lualatex+MiKTeX leaves these black, which
% looks inconsistent with the rest of the blue accent scheme.
\renewcommand*{\firstnamestyle}[1]{{\fontsize{34}{36}\bfseries\upshape\color{color1}#1}}
\renewcommand*{\lastnamestyle}[1]{{\fontsize{34}{36}\bfseries\upshape\color{color1}#1}}
\renewcommand*{\sectionstyle}[1]{{\sectionfont\color{color1}#1}}

\usepackage[utf8]{inputenc}
\usepackage{hyperref}
\hypersetup{
    colorlinks=true,
    linkcolor=blue,
    filecolor=magenta,
    urlcolor=blue,
    pdftitle={Andrea Piccolo - CV},
    pdfpagemode=FullScreen,
}
\usepackage[scale=0.77]{geometry}
\usepackage{needspace}
\usepackage{import}
\hyphenation{FEniCSx StonedFEniCSx}

% personal data
\name{Andrea}{Piccolo}
\address{Leeds, UK}{}{}
\phone[mobile]{+44 7831 007579}
\email{andreapiccolope@gmail.com}
\extrainfo{\href{https://www.linkedin.com/in/ap0889/}{LinkedIn}, \href{https://github.com/APiccolo89}{GitHub} \textbar{} UK Global Talent Visa -- no sponsorship required}

\begin{document}

\makecvtitle

% ============================================================
%     PROFILE STATEMENT
% ============================================================

\vspace{6pt}
\small{Python software engineer with a PhD in computational geodynamics and ten years building numerical modelling software. I own code end to end: design, implementation, automated tests, CI, documentation, and the workflows that run it. Turned a single-script model into a tested, object-oriented open-source package. Use Claude Code daily for writing, testing, and reviewing code.}

% ============================================================
%     SKILLS
% ============================================================

\section{Skills}
\vspace{1pt}
\begin{itemize}

\item \textbf{Python Software Engineering}: Object-oriented design (class hierarchies, dataclass configuration, YAML input), packaging (pyproject, pip), NumPy, SciPy, pandas, h5py, Numba; also Julia, C, MATLAB

\item \textbf{Testing \& Code Quality}: pytest suites with Codecov coverage, regression tests against published benchmarks, GitHub Actions CI, Sphinx documentation, contributions to shared open-source codebases

\item \textbf{Modelling \& Optimisation}: Numerical models (finite element, finite difference); reduced-order models; optimisation algorithm for parameter recovery against experimental data

\item \textbf{Data Pipelines \& Compute}: Automated simulation and post-processing workflows; TB-scale data processing (20 TB condensed to 200 GB public databases); HPC clusters (SLURM, PBS), MPI; Linux; SQL (learning)

\item \textbf{AI-Assisted Development}: Daily use of Claude Code as an agentic coding assistant for development, test writing, debugging, and code review

\end{itemize}

% ============================================================
%     PROFESSIONAL EXPERIENCE
% ============================================================

\section{Professional Experience}
\vspace{3pt}
\begin{itemize}

\needspace{5\baselineskip}
\item{\cventry{02/2025--Present}{Research Fellow}{University of Leeds}{Leeds, UK}{}{\vspace{1pt}
\begin{itemize}
    \item Built \href{https://github.com/APiccolo89/StonedFenicsx}{StonedFEniCSx}, an open-source FEniCSx package for coupled thermal and flow modelling, growing it from a single script into an object-oriented package: problem and solver class hierarchies, dataclass configuration, YAML input, pip-installable layout, so new models are set up without touching the solver
    \item Set up CI with GitHub Actions: pytest suite with Codecov coverage, regression tests against published benchmarks, and Sphinx documentation deployed to GitHub Pages
    \item Built automated workflows to run simulations, post-process outputs, and analyse large datasets
\end{itemize}}}

\vspace{3pt}

\needspace{5\baselineskip}
\item{\cventry{11/2021--05/2024}{Research Fellow}{University of Bayreuth}{Bayreuth, Germany}{}{\vspace{1pt}
\begin{itemize}
    \item Designed a 0D reduced-order model predicting full 3D simulation behaviour, and an optimisation algorithm that recovers the parameter combinations matching 2D experimental data
    \item Delivered public databases used by other groups: 20 TB of simulated data condensed into 200 GB of accessible outputs
    \item Extended an open-source Julia package, coupling forward mechanical modelling with basin modelling
    \item Wrote Python tools for data cleaning, analysis, and visualisation used in project reporting; supplied data and support to the KIT Karlsruhe modelling group
\end{itemize}}}

\vspace{3pt}

\needspace{5\baselineskip}
\item{\cventry{11/2015--07/2021}{Research Fellow}{Johannes Gutenberg Universität Mainz}{Mainz, Germany}{}{\vspace{1pt}
\begin{itemize}
    \item Incorporated chemical modelling into a large-scale 3D numerical code (\href{https://github.com/UniMainzGeo/LaMEM}{LaMEM}) developed and used by several groups
    \item Generated and processed 100+ TB of simulation data; developed 2D/3D visualisation tools to identify physical trends
    \item Presented results at international conferences; supported supervision and teaching
\end{itemize}}}

\end{itemize}

% ============================================================
%     EDUCATION
% ============================================================

\needspace{8\baselineskip}
\section{Education}
\vspace{1pt}
\begin{itemize}

\item{\cventry{11/2015--05/2019}{PhD in Computational Geodynamics}{Johannes Gutenberg Universität Mainz}{Mainz, Germany}{}{\vspace{1pt}
Thesis: ``Continental crust production during the early Earth: Constraints from geodynamic models.'' Thermo-mechanical numerical modelling, data analysis, CI.
}}

\vspace{3pt}

\item{\cventry{10/2012--03/2015}{MSc in Geology and Technical Geology}{University of Padua}{Padua, Italy}{}{\vspace{1pt}
Data analysis, numerical modelling, scientific communication.
}}

\vspace{3pt}

\item{\cventry{10/2008--03/2012}{BSc in Earth Sciences}{University of Padua}{Padua, Italy}{}{\vspace{1pt}
Data analysis, MATLAB.
}}

\end{itemize}

% ============================================================
%     REPOSITORIES & SELECTED PUBLICATIONS
% ============================================================

\section{Repositories}
\vspace{1pt}
\begin{itemize}
\item \textbf{Author:} \href{https://github.com/APiccolo89/StonedFenicsx}{StonedFEniCSx} \textbar{} \textbf{Contributor:} \href{https://github.com/UniMainzGeo/LaMEM}{LaMEM}, \href{https://github.com/JuliaGeodynamics/GeophysicalModelGenerator.jl}{GeophysicalModelGenerator.jl}, \href{https://bitbucket.org/bkaus/mvep2/}{MVEP2}
\end{itemize}

\needspace{5\baselineskip}
\section{Selected Publications}
\vspace{1pt}
{\small
\begin{itemize}
\item Kaus, B.\,J., \ldots, \textbf{Piccolo, A.}, et al.\ (2024). GeophysicalModelGenerator.jl: A Julia package to visualise geoscientific data and create numerical model setups. \textit{Journal of Open Source Software}, 9(103), 6763.
\item \textbf{Piccolo, A.}, Spang, A., et al.\ (2026). The dynamics and surface signal of slab break-off in continental settings: Insights from 3D numerical modeling. \textit{JGR: Solid Earth}, 131(5).
\item \textbf{Piccolo, A.}, Palin, R.\,M., et al.\ (2019). Generation of Earth's early continents from a relatively cool Archean mantle. \textit{Geochem., Geophys., Geosyst.}, 20(4), 1679--1697.
\item[] 12 peer-reviewed papers in total (4 as first author); full list on request.
\end{itemize}}

% ============================================================
%     LANGUAGES
% ============================================================

\section{Languages}
\vspace{1pt}
\begin{itemize}
\item Italian (native), English (C1), German (B1).
\end{itemize}

% ============================================================
%     REFERENCES
% ============================================================

\section{References}
\vspace{1pt}
\begin{itemize}
\item Prof.\ Dr.\ Timothy J.\ Craig --- University of Leeds (\href{mailto:T.J.Craig@leeds.ac.uk}{T.J.Craig@leeds.ac.uk})
\item Prof.\ Dr.\ Marcell Thielmann --- University of Bonn (\href{mailto:thielmann@geo.uni-bonn.de}{thielmann@geo.uni-bonn.de})
\item Prof.\ Dr.\ Nevena Andri\'c Toma\v{s}evi\'c --- KIT Karlsruhe (\href{mailto:nevena.tomasevic@kit.edu}{nevena.tomasevic@kit.edu})
\item Dr.\ Arne Spang --- University of Bayreuth (\href{mailto:arne.spang@uni-bayreuth.de}{arne.spang@uni-bayreuth.de})
\end{itemize}

\end{document}
